%% ============================================================
%%  Chapter 1: Introduction
%% ============================================================
\chapter{Introduction}
\label{chap:introduction}


%% ── 1.1 ─────────────────────────────────────────────────────────
\section{Motivation}
\label{sec:intro_motivation}

Modern warehouses operate at a scale that exceeds the reach of any centralized
controller.
Amazon's fulfillment centres deploy tens of thousands of drive units; Ocado's
customer fulfillment centres run more than a thousand robots on a single grid
floor.
At peak throughput, each robot receives a new task within seconds of completing
the last one, and a fleet of hundreds must navigate shared aisles without
collision, without deadlock, and without leaving any induct station backlogged.
The coordination problem is, at its core, a continuous assignment problem:
which robot should pick up which item, in what order, along which path, and
when should the decision be revised?

Two broad strategies have emerged for tackling this problem.
\emph{Centralized} approaches gather all task and robot state at a single
planner, compute a globally optimal assignment, and broadcast instructions to
the fleet.
These methods enjoy strong theoretical guarantees. The Sequential Greedy
Algorithm~\cite{minoux1978accelerated} achieves at least half the optimal
assignment value in polynomial time and their simulation-to-deployment gap
is well understood.
Their weaknesses, however, are equally well documented: a single coordinator
represents a point of failure, scaling the coordinator's computation with fleet
size is non-trivial, and the assumption of reliable broadcast communication
is rarely met in real warehouses, where range-limited radios, metal shelving,
and high robot density produce an intermittently connected network of links.

\emph{Decentralized} approaches distribute the coordination decision across
the fleet.
Each robot runs a local planning algorithm using only the information available
within its communication neighbourhood, achieving resilience to node failures
and communication outages at the cost of solution quality guarantees.
The Consensus-Based Bundle Algorithm (CBBA)~\cite{choi2009consensus} and its
global extension GCBBA~\cite{kha2025enhancing} represent the state of the art
in this family: they provably converge to within 50\% of the optimal
assignment value on any \emph{connected} communication graph.

The critical qualification of \emph{connected} is where both methods fall
short of real-world requirements.
In a warehouse with a physical communication range of 5--10 meters,
robots spread across a 30-meter floor will frequently be in communication with
only a subset of the fleet.
When the communication graph fragments, CBBA and GCBBA may converge to
conflicting assignments (multiple robots claiming the same task) or leave
isolated robots idle.
Neither method specifies a recovery protocol for reconnection events or for
tasks that arrive while components are disconnected.
The result is a performance cliff: as the communication range drops below
the warehouse diameter, allocation quality degrades sharply or fails entirely.

Beyond connectivity, real warehouse deployments impose a second class of
constraints that academic benchmarks routinely abstract away:
\emph{operational constraints} arising from physical robot limitations.
Drive units carry finite batteries; a robot that exhausts its charge
mid-aisle becomes a stationary obstacle requiring manual retrieval.
Idle robots parked in main picking aisles force every active robot to
route around them, degrading throughput even when the fleet is not at
capacity.
These constraints interact with allocation in non-trivial ways: a task that
appears attractive to a low-energy agent may be unreachable before the
agent must divert to a charger; an idle agent holding position in a busy
corridor may block the paths of several active agents simultaneously.
A coordination system evaluated only on path quality and task completion
rate --- without modelling energy or congestion from idle robots --- will
miss these failure modes entirely and may perform poorly when deployed
on real hardware.

This thesis develops the \textbf{Local Consensus-Based Bundle Algorithm
(LCBA)}, which removes the connectivity assumption from GCBBA while preserving
its convergence guarantees within each connected component and extending
its correctness to dynamic, intermittently disconnected topologies.
LCBA is integrated with a family of three collision-free path planning strategies
to form a complete multi-robot coordination pipeline that can be evaluated
across a wide range of operating conditions.

The practical scope extends beyond warehouse automation.
Underwater robotic surveys operate without GPS and with highly attenuated
acoustic communication, producing exactly the intermittent connectivity
structure that LCBA is designed to handle.
Disaster-response scenarios where physical infrastructure is damaged and
communication links are unreliable require decentralized coordination
methods that continue to function even when large portions of the team are
out of contact.
Space-based robotic systems face communication delays and blackout periods
that similarly preclude any centralized coordinator.
While this thesis focuses on the warehouse setting (which provides a reliable
set of measurable metrics and a reproducible simulation environment), the
algorithm is designed with these broader application domains in mind.


%% ── 1.2 ─────────────────────────────────────────────────────────
\section{Problem Statement}
\label{sec:intro_problem}

This thesis addresses two coupled sub-problems in a $30 \times 30$ gridworld
warehouse environment populated by $\Na = 6$ mobile agents,
$|\mathcal{S}^{\text{in}}| = 8$ induct stations, and
$|\mathcal{S}^{\text{ej}}| = 38$ eject stations.

\textbf{Sub-problem 1: Decentralized multi-robot task allocation under dynamic,
intermittently disconnected communication graphs.}
Tasks arrive continuously at induct stations over time, not as a fixed batch at
$t = 0$.
Agents communicate over a proximity-based wireless channel whose topology
changes every timestep as agents move.
The communication graph may be arbitrarily disconnected: in the worst case,
each agent forms its own isolated component.
The allocation method must correctly assign tasks, resolve conflicts, and
reassign work as the topology changes. More importantly, it must operate
using only information locally
available within each agent's communication neighbourhood.
Existing methods (CBBA, GCBBA, DMCHBA) break down or stall when the graph is
persistently disconnected.

\textbf{Sub-problem 2: Collision-free execution of allocated tasks.}
Once tasks are assigned, agents must navigate from their current positions to
the assigned induct station, pick up the item, and deliver it to the assigned
eject station, all without occupying the same grid cell as another agent at the
same timestep.
The path planning method must be computationally feasible within a per-timestep
budget and must remain functional as the number of agents or the conflict
density increases.

The two sub-problems are coupled: travel-time estimates used in the allocation
algorithm depend on the paths that agents will follow, which in turn depend on
the decisions of other agents.
A practical system must handle this coupling through event-driven replanning
rather than by solving the joint problem exactly.


%% ── 1.3 ─────────────────────────────────────────────────────────
\section{Research Questions}
\label{sec:intro_questions}

The thesis investigates three research questions:

\begin{description}
  \item[\textbf{RQ1.}] Can LCBA maintain task throughput under communication
    conditions where static allocation fails?
    Specifically, does event-driven replanning allow LCBA to complete tasks
    even when the communication graph is persistently disconnected, and by how
    much does it outperform CBBA, SGA, and DMCHBA at low communication ranges?

  \item[\textbf{RQ2.}] How close to optimal (SGA) does LCBA perform as
    communication connectivity varies?
    On connected graphs where SGA is a valid reference, what is the
    sub-optimality ratio of LCBA, and does it remain within the theoretical
    50\%-optimality bound?

  \item[\textbf{RQ3.}] How do alternative path planning strategies affect
    system throughput, task wait time, and computational overhead?
    Is there a single planner that dominates across all operating conditions,
    or does the optimal choice depend on task load and agent count?
\end{description}


%% ── 1.4 ─────────────────────────────────────────────────────────
\section{Contributions}
\label{sec:intro_contributions}

This thesis makes the following contributions:

\begin{enumerate}
  \item \textbf{The Local Consensus-Based Bundle Algorithm (LCBA).}
        LCBA extends GCBBA to handle three challenges that GCBBA and CBBA
        cannot:
        (a) dynamic, online task arrivals via event-driven replanning;
        (b) arbitrarily disconnected communication graphs by restricting
            consensus to each connected component and resolving conflicts
            upon reconnection; and
        (c) dynamically changing communication topology by reconstructing
            the graph at every timestep.
        On connected components, LCBA inherits GCBBA's convergence to the SGA
        solution and the associated 50\%-optimality guarantee.

  \item \textbf{A hybrid, modular coordination pipeline.}
        LCBA is integrated with three interchangeable path planning strategies
        (Cooperative A*, Priority-Based Search, and Rolling Horizon Collision
        Resolution) within a single orchestration framework.
        The modular design allows the allocation and path planning components
        to be studied independently.

  \item \textbf{A dual-mode warehouse simulation.}
        The simulation supports both a \emph{batch mode} (fixed task set,
        makespan objective) for comparison with prior work, and a
        \emph{steady-state mode} (continuous task arrivals, throughput
        objective under fixed simulation and compute budgets) that better
        reflects the infinite-horizon operation of real warehouse systems.
        Both modes share the same environment, agents, and infrastructure,
        enabling direct comparison.

  \item \textbf{Empirical evaluation across 432+ experimental runs.}
        The system is evaluated across a full factorial sweep of communication
        range (6 values from fully isolated to fully connected),
        task arrival rate (12 values from light load to extreme saturation),
        path planner (3 variants), and allocation method (LCBA, CBBA, SGA, DMCHBA),
        across six environments ranging from 6 to 200 agents.
        This scale of evaluation provides statistical confidence in the results
        and reveals the interaction effects between communication conditions
        and planning strategy.

  \item \textbf{Operationally realistic simulation with charging, wait stations,
        and role-separated planners.}
        The simulation incorporates finite agent energy with proactive charging
        (agents navigate to charging stations before running out of power, with an
        energy-budget filter that prevents bids on unaffordable tasks),
        designated wait stations where idle agents park to avoid obstructing active
        corridors, and independently configurable path planners for task-executing,
        charging, and idle-parking agent roles.
        Together these features bridge the gap between typical academic MAPF benchmarks
        and the operational constraints of real warehouse deployments, ensuring that
        LCBA is evaluated under realistic pressure and that its design choices are
        exercised by the conditions they were designed to handle.
\end{enumerate}


%% ── 1.5 ─────────────────────────────────────────────────────────
\section{Thesis Outline}
\label{sec:intro_outline}

The remainder of this thesis is organized as follows.
\Cref{chap:background} reviews related work in multi-robot task allocation,
multi-agent path finding, communication-constrained coordination, and warehouse
automation, and identifies the specific gap this thesis fills.
\Cref{chap:methodology} describes the system architecture, the LCBA algorithm
and its warehouse adaptation, the three path planning strategies, the dynamic
replanning mechanism, and the baseline methods.
\Cref{chap:experiments} presents the experimental setup (including the dual
batch/steady-state design), defines all metrics, and reports quantitative
results across the full parameter sweep.
\Cref{chap:conclusion} summarizes the findings, revisits the contributions
in light of the experimental evidence, discusses limitations, and outlines
directions for future work.